\subsection{Tangent lines, normal lines, and units}

If \(f\) is differentiable at \(x=a\), the tangent line is
\[
y-f(a)=f'(a)(x-a).
\]
When \(f'(a)\ne0\), the normal line has slope \(-1/f'(a)\).

\begin{examplebox}[Tangent and normal to a parabola]
For \(f(x)=x^2\) at \(a=1\),
\[
f(1)=1,
\qquad
f'(1)=2.
\]
Hence the tangent is
\[
y-1=2(x-1),
\]
and the normal is
\[
y-1=-\frac12(x-1).
\]
\end{examplebox}

\begin{interpretationbox}[Units of a derivative]
If \(x\) is measured in seconds and \(f(x)\) in metres, then \(f'(x)\) is measured in metres per second. A derivative always has units of output divided by units of input. This is essential in applications: it distinguishes position, velocity, acceleration, concentration change, and marginal cost.
\end{interpretationbox}

\begin{remarkbox}
A numerical derivative without units may be mathematically correct but incomplete as an applied conclusion.
\end{remarkbox}
